\chapter{LITEN, DTS et INES}\label{ch:liten-dts-et-ines}

L'activité du \gls{LITEN} est découpée en 5 parties dans son dernier rapport d'activités\cite{CEA_prez_LITEN} :
\begin{enumerate}
    \item Photovoltaïque
    \item Batterie
    \item Hydrogène, molécules et carburants durables
    \item Matériaux et économie circulaire
    \item Systèmes, réseaux et efficacité énergétique
\end{enumerate}

Quelques chiffres viennent quantifier cette activité :
\begin{itemize}
    \item 1000 chercheurs, techniciens et membres des équipes de support
    \item 200 publications annuelles
    \item 200 doctorants et post-doc
    \item 160 millions d'euros de budget annuel
    \item 200 partenaires industriels
    \item 12 plateformes technologiques
    \item 2000 familles de brevets en portefeuille
\end{itemize}

\begin{figure}[H]
    \centering
    \fbox{\scalebox{0.9}{\begin{tikzpicture}[
    main/.style={text centered, text width=8cm, fill=blue!20, draw, rectangle, rounded corners, font=\bfseries\Large},
    département/.style={text centered, text width=3cm, fill=gray!50, draw, rectangle, rounded corners},
    service/.style={grow=down, xshift=4cm, text centered, text width=5cm, draw, rectangle, rounded corners, edge from parent path={(\tikzparentnode.225) |- (\tikzchildnode.west)}},
    first/.style    ={level distance=15ex},
    second/.style   ={level distance=30ex},
    third/.style    ={level distance=45ex},
    fourth/.style   ={level distance=60ex},
    fifth/.style    ={level distance=75ex},
    level 1/.style={sibling distance=12.5em, level distance=8em}]

    \node[main, anchor=south]{Direction du LITEN}
    [edge from parent fork down]
    %
    child{node (d1) [département] {\textbf Département \\ des \\ \textbf Technologies \textbf Solaires}
        child[service,first]  {node {\textbf Service Procédés \\ \textbf Cellules \textbf Photo\textbf{v}oltaïques}}
        child[service,second] {node {\textbf Service \textbf Modules \\et \textbf Systèmes \textbf Photovoltaïques}}
        child[service,third]  {node {\textbf Service \textbf Intégration et \textbf Réseaux \textbf Énergétiques}}
    }
    %
    child{node (d2) [département] {\textbf Département \\ \textbf Électricité \\ et \textbf Hydrogène pour les \textbf Transports}}
    %
    child{node (d3) [département] {\textbf Département \\ \textbf Thermique \\ \textbf Conversion et \textbf Hydrogène}}
    %
    child{node (d4) [département] {\textbf Département \\ \textbf Technologies \\ des \textbf Nouveaux \textbf Matériaux}};

\end{tikzpicture}}}
    \caption{Organigramme simplifié du LITEN}
    \label{fig:organigramme-LITEN}
\end{figure}

Au sein du \gls{LITEN}, le \gls{DTS} a pour mission de développer des composants solaires et de les intégrer dans des structures et infrastructures de \("\)PV Everywhere\("\) et plus largement dans des systèmes électriques en y associant l'intelligence artificielle pour faciliter leur insertion.
Ses installations techniques sont situées au sein de l'\gls{INES} sur le campus Technolac près du lac du Bourget.

Le site est créé en 2005 sous l'impulsion d'élus locaux, d'industriels historiques de la région et des organismes de recherche \textit{(CEA, CNRS et Université de Savoie)}.
Ce lancement fait suite à une discussion de plusieurs années dont sont notamment à l'initiative Michel Barnier et Jean-Pierre Vial.

L'\gls{INES} est composé pour la grande majorité d'acteurs du \gls{CEA}, complétés par plusieurs start-ups, et quelques organismes de formation.
Il poursuit ainsi le triple but de mener la R\&D pour les technologies solaires avancées, d'accompagner les territoires et les entreprises à renforcer leurs capacités et connaissances en énergie solaire, et de représenter la filière française dans différentes instances régionales, nationales, européennes et mondiales\cite{site_ines}.

\begin{bluebox}
    Interrogé sur les enjeux pour le photovoltaïque en France et sur la stratégie du \gls{LITEN} pour y répondre\cite{CEA_prez_LITEN}, le chef du \gls{DTS} mentionne le besoin d'intégration verticale de la technologie PV en Europe, c'est-à-dire la maîtrise de l'ensemble de la chaîne de valeur,  \("\)depuis le polysilicium jusqu'au module\("\).
    La R\&D du département est à ce titre focalisée sur l'optimisation des procédés de fabrication des technologies actuelles comme la réduction des temps de passage de chaque procédé.
    Un autre axe de travail du \gls{DTS} est l'intégration des panneaux PV sur des surfaces déjà artificialisées, l'intégration de l'énergie solaire en réseau, notamment par le passage en courant continu pour diminuer les pertes électriques.
    Le dernier point soulevé est celui de soutenabilité des technologies solaires en termes de consommation de matières premières, et l'acceptabilité sociétale des technologies proposées.

    Cette stratégie est porteuse de résultats, que l'on peut retrouver dans le rapport d'activités 2025 du \gls{CEA}~\cite{CEA_prez_generale} où est présentée la série de records de rendements atteints par les cellules PV développées en partenariat avec 3SUN\@.
    Le dernier en date est de 30,8\%, un chiffre dépassant les limites de rendement du photovoltaïque conventionnel au silicium.
\end{bluebox}

